% -*- coding: utf-8 -*-
% Preamble for md-to-tex-slides.md (included via `include-in-header` in
% md-to-tex-beamer.yaml).  It lives in a file because pandoc parses
% `header-includes` values as Markdown, which mangles raw LaTeX such as
% fontspec's `[RawFeature={fallback=...}]`.
\usepackage{fontspec}
\directlua{luaotfload.add_fallback("emojifb", {"Segoe UI Emoji:mode=harf"})}
\setsansfont{Latin Modern Sans}[RawFeature={fallback=emojifb}]
\setmainfont{Latin Modern Roman}[RawFeature={fallback=emojifb}]
\setmonofont{Latin Modern Mono}[RawFeature={fallback=emojifb},Scale=0.86]

\usepackage{amsmath,amssymb}
\usepackage{xcolor}
\definecolor{nordred}{HTML}{BF616A}
\definecolor{nordblue}{HTML}{5E81AC}
\definecolor{nordgreen}{HTML}{A3BE8C}
\definecolor{nordyellow}{HTML}{EBCB8B}
\definecolor{nordpurple}{HTML}{B48EAD}
\definecolor{norddark}{HTML}{2E3440}
\definecolor{nordlight}{HTML}{ECEFF4}

\usepackage{pgf,tikz}
\usetikzlibrary{arrows.meta,positioning,shapes.geometric,calc,fit,backgrounds}
\tikzset{
  nb/.style={draw, rounded corners, align=center, font=\scriptsize,
             inner sep=4pt, minimum height=7mm},
  nblue/.style={nb, fill=blue!12,  draw=nordblue,  line width=1pt},
  nyellow/.style={nb, fill=yellow!35, draw=orange!80!black, line width=1pt},
  npurple/.style={nb, fill=violet!14, draw=nordpurple, line width=1pt},
  nred/.style={nb, fill=red!12, draw=nordred, line width=1.5pt},
  ngreen/.style={nb, fill=green!16, draw=green!45!black, line width=1pt},
  npink/.style={nb, fill=magenta!10, draw=magenta!70!black, line width=1pt},
  ar/.style={-{Stealth[length=2mm]}, draw=black!55, line width=0.8pt},
}

\usepackage{listings}
\lstset{
  basicstyle=\ttfamily\scriptsize,
  breaklines=true,
  showstringspaces=false,
  columns=fullflexible,
  frame=single,
  rulecolor=\color{black!20},
}

\usetheme{Frankfurt}
\setbeamertemplate{navigation symbols}{}
\setbeamercolor{structure}{fg=nordblue}
\setbeamercolor{frametitle}{fg=norddark,bg=nordlight}
\setbeamercolor{math text}{fg=nordblue}
\setbeamercovered{transparent}

\newcommand{\columnsbegin}{\begin{columns}}
\newcommand{\columnsend}{\end{columns}}
\newcommand{\col}[1]{\column{#1}}
